\section{Las capacidades fundamentales que no cambian}

\subsection{Habilidades que la IA no puede sustituir}

Aunque las herramientas y los flujos de trabajo cambian con rapidez, la idea central que Mihail Eric planteó en la semana 1 sigue siendo válida para la próxima década:

\begin{importantbox}{Capacidades fundamentales que atraviesan los ciclos tecnológicos}
\begin{enumerate}
    \item \textbf{Comprensión del negocio}: entender las necesidades del usuario y el valor comercial; la IA no sabe «qué hay que hacer»
    \item \textbf{Pensamiento arquitectónico}: las decisiones de diseño a nivel de sistema exigen sopesar innumerables trade-offs
    \item \textbf{Buen criterio (Taste)}: la capacidad de distinguir entre código bueno y malo, entre buen diseño y mal diseño
    \item \textbf{Capacidad de lectura de código}: revisar y entender el código generado por la IA importa más que escribirlo uno mismo
    \item \textbf{Comunicación y coordinación}: ya sea que se gestionen personas o Agentes, la comunicación clara es la base
\end{enumerate}
\end{importantbox}

\subsection{Recomendaciones para la nueva generación de desarrolladores}

Con base en los aprendizajes de todo el curso, algunas recomendaciones para los desarrolladores del futuro:

\begin{enumerate}
    \item \textbf{Aprende las herramientas de IA, pero no dependas de ellas}: entender los principios subyacentes (semanas 1--2) perdura más que dominar una herramienta en particular
    \item \textbf{Invierte en la capacidad de leer código}: en la era de la programación con IA, la capacidad de revisión vale más que la de escritura (semana 7)
    \item \textbf{Practica la delegación y el trabajo en múltiples hilos}: es la habilidad central del patrón Agent Manager (semanas 3--4)
    \item \textbf{Entiende los principios básicos de seguridad}: la IA introduce nuevas superficies de ataque; la alfabetización en seguridad no es negociable (semana 6)
    \item \textbf{Adopta la experimentación rápida}: como dijo Mihail en la semana 1, «no hay un paradigma maduro, todos están explorando a tientas»
    \item \textbf{Construye para el modelo de dentro de 6 meses}: la recomendación de Boris Cherny (semana 4) aplica a la elección de todas las herramientas y flujos de trabajo
\end{enumerate}

\subsection{La evolución de la forma de los equipos futuros}

Si se traducen las conclusiones de todo el curso al plano organizacional, es muy probable que los equipos de software del futuro presenten una nueva estratificación:

\begin{table}[H]
\centering
\begin{tabular}{p{0.2\textwidth} p{0.26\textwidth} p{0.28\textwidth}}
\toprule
Nivel del equipo & Forma principal de trabajo & Origen de la ventaja competitiva \\
\midrule
Nivel de desarrollador individual & Completar tareas de extremo a extremo con IDEs de IA y Agentes & Velocidad de prueba y error, capacidad de ejecución interfuncional \\
Nivel de equipo pequeño & Varias personas gestionan cada una varios Agentes y los integran de forma colaborativa & Capacidad de planeación, límites de las interfaces y mecanismos de revisión \\
Nivel de equipo de plataforma & Mantener una cadena de herramientas, normas, evaluaciones y una línea base de seguridad unificadas & Capacidad de replicar a escala organizacional \\
\bottomrule
\end{tabular}
\caption{Cambios en la forma de los equipos en la era de la ingeniería de software con IA}
\end{table}

\begin{knowledgebox}{La ventaja organizacional futura vendrá más del diseño de sistemas que de expertos individuales}
Cuando la eficiencia de codificación individual se ve amplificada enormemente por la IA, lo que suele marcar la diferencia real ya no es «quién escribe más rápido», sino «quién logra ensamblar herramientas, procesos, seguridad y sistemas de revisión en un sistema estable». Por eso todo el curso regresa, al final, a la perspectiva de la organización y los procesos.
\end{knowledgebox}

\subsection{Resumen de la sección}

Cuanto más rápido cambia la tecnología, más importantes son las capacidades fundamentales. La comprensión del negocio, el pensamiento arquitectónico, el buen criterio, la lectura de código y la comunicación son las habilidades permanentes que atraviesan el ciclo tecnológico de la IA.

